\documentclass[sigconf]{acmart}

\usepackage{booktabs}
\usepackage{graphicx}
\usepackage{listings}
\usepackage{xcolor}

\AtBeginDocument{%
  \providecommand\BibTeX{{Bib\TeX}}
}

\setcopyright{none}
\acmConference[SABD 2025/26]{Sistemi e Architetture per Big Data}{Giugno 2026}{Roma, Italia}

\begin{document}

\title{Analisi di Dati sui Voli con Apache Spark}
\subtitle{Progetto 1 — Sistemi e Architetture per Big Data, A.A. 2025/26}

\author{Luca Carloni}
\affiliation{%
  \institution{Università degli Studi di Roma ``Tor Vergata''}
  \department{Dipartimento di Ingegneria Civile e Informatica}
}
\email{TODO@students.uniroma2.eu}

% Se il gruppo è composto da più studenti, aggiungere gli altri autori qui
% \author{Nome Cognome}
% \affiliation{...}
% \email{...}

\begin{abstract}
In questo lavoro presentiamo un sistema Big Data per l'analisi di dati storici
sui voli aerei statunitensi, relativi al periodo gennaio--aprile 2025 (circa
2.200.000 eventi). L'architettura adottata integra Apache NiFi per la data
ingestion, HDFS per lo storage distribuito, Apache Spark per il processing
(DataFrame API e Spark SQL), Redis per l'export dei risultati e Grafana per la
visualizzazione. Vengono implementate tre query analitiche sui ritardi e sui
tassi di cancellazione, un'analisi di clustering K-Means sulle compagnie aeree
e un confronto sperimentale tra DataFrame API e Spark SQL in termini di tempi
di processamento.
\end{abstract}

\maketitle

% ─────────────────────────────────────────────────────────────────────────────
\section{Introduzione}

L'analisi di grandi volumi di dati sui trasporti aerei è un caso d'uso
emblematico per le piattaforme Big Data: i dati sono strutturati, numerosi e
soggetti a query di aggregazione complesse. In questo progetto utilizziamo il
dataset fornito dal Dipartimento dei Trasporti degli Stati Uniti, contenente
informazioni su tutti i voli operati nel periodo 1 gennaio--30 aprile 2025.

L'obiettivo è rispondere a tre query analitiche sui ritardi, calcolare il tasso
di cancellazione e applicare tecniche di machine learning (clustering K-Means)
per identificare gruppi di compagnie aeree con comportamento operativo simile.
L'intero sistema è deployato via Docker Compose su un nodo standalone.

% ─────────────────────────────────────────────────────────────────────────────
\section{Architettura del Sistema}

\begin{figure}[h]
  \centering
  % Inserire qui lo schema architetturale (immagine o TikZ)
  % \includegraphics[width=\linewidth]{architecture.png}
  \fbox{\parbox{0.9\linewidth}{\centering\vspace{1.5cm}
    [Schema architetturale da inserire]\vspace{1.5cm}}}
  \caption{Architettura del sistema: NiFi → HDFS → Spark → Redis → Grafana.}
  \label{fig:arch}
\end{figure}

L'architettura, illustrata in Figura~\ref{fig:arch}, segue un pipeline
lineare composto da cinque strati:

\begin{enumerate}
  \item \textbf{Data Ingestion (Apache NiFi 1.25.0)}: un flusso NiFi legge i
        CSV grezzi dal filesystem locale, li converte in formato Parquet
        (tramite \texttt{ConvertRecord} con schema Avro) e li carica su HDFS.

  \item \textbf{Storage distribuito (Apache Hadoop HDFS 3.2.1)}: i file
        Parquet vengono memorizzati in \texttt{/sabd/processed/} con fattore
        di replica 1 (nodo singolo). I risultati delle query vengono scritti
        in \texttt{/sabd/results/}.

  \item \textbf{Processing (Apache Spark 3.5.0)}: i job Python (PySpark)
        leggono i Parquet da HDFS, eseguono le query mediante DataFrame API o
        Spark SQL e producono i risultati in CSV (locale e HDFS).

  \item \textbf{Export (Redis 7.2)}: i risultati vengono caricati in Redis
        tramite strutture HASH e ZSET, organizzate per essere direttamente
        interrogate da Grafana.

  \item \textbf{Visualizzazione (Grafana 10.4.2)}: quattro dashboard
        pre-configurate (Q1, Q2, Q3, Clustering) visualizzano i risultati
        attraverso il plugin \texttt{redis-datasource}.
\end{enumerate}

Tutti i servizi comunicano su una rete Docker bridge (\texttt{sabd-net}) e
sono orchestrati tramite un singolo file \texttt{docker-compose.yml}.

% ─────────────────────────────────────────────────────────────────────────────
\section{Dataset e Data Ingestion}

\subsection{Dataset}

Il dataset contiene 2.200.000 eventi che descrivono i voli effettuati negli
Stati Uniti da gennaio ad aprile 2025, distribuiti in quattro file CSV mensili.
Per gli scopi del progetto sono stati selezionati 14 campi: identificativo
compagnia (\texttt{OP\_UNIQUE\_CARRIER}), data (\texttt{YEAR}, \texttt{MONTH},
\texttt{DAY\_OF\_MONTH}), orario schedulato di partenza
(\texttt{CRS\_DEP\_TIME}), ritardi (\texttt{DEP\_DELAY}, \texttt{ARR\_DELAY}),
stato del volo (\texttt{CANCELLED}, \texttt{DIVERTED}) e le cinque componenti
di ritardo (\texttt{CARRIER\_DELAY}, \texttt{WEATHER\_DELAY},
\texttt{NAS\_DELAY}, \texttt{SECURITY\_DELAY},
\texttt{LATE\_AIRCRAFT\_DELAY}).

\subsection{Pipeline di Ingestion}

Il flusso NiFi adotta la seguente catena di processori:

\begin{enumerate}
  \item \textbf{ListFile}: scansiona ricorsivamente la directory
        \texttt{/opt/nifi/data/raw} ogni 10 secondi, individuando i file CSV
        mensili tramite pattern regex.
  \item \textbf{FetchFile}: legge il contenuto del file.
  \item \textbf{ConvertRecord}: converte ogni record CSV in formato Parquet
        applicando uno schema Avro che tipizza correttamente i campi
        numerici (\texttt{int} per date/orari, \texttt{double} per ritardi).
  \item \textbf{UpdateAttribute}: rinomina l'estensione da \texttt{.csv}
        a \texttt{.parquet}.
  \item \textbf{PutHDFS}: scrive il file in \texttt{/sabd/processed/} su HDFS,
        con strategia di conflitto \texttt{replace}.
\end{enumerate}

La scelta del formato Parquet, con granularità mensile, permette a Spark di
sfruttare la predicate pushdown e la codifica a colonne, riducendo l'I/O
durante le query di aggregazione.

% ─────────────────────────────────────────────────────────────────────────────
\section{Query e Risultati}

\subsection{Q1 — Statistiche mensili AA vs DL}

La query Q1 calcola, per ciascun mese del periodo in esame, le statistiche di
\texttt{DEP\_DELAY} (media, minimo, massimo) e il tasso di cancellazione per
le compagnie American Airlines (AA) e Delta (DL). Il calcolo dei ritardi
considera esclusivamente i voli non cancellati (\texttt{CANCELLED = 0}),
mentre il cancellation rate è calcolato su tutti i voli del gruppo.

I risultati sono riportati nella Tabella~\ref{tab:q1}.

\begin{table}[h]
\caption{Q1 — Statistiche mensili AA e DL (gen--apr 2025).}
\label{tab:q1}
\small
\begin{tabular}{llrrrr}
\toprule
Carrier & Mese & Avg Delay & Min & Max & Canc. Rate \\
        &      & (min)     & (min) & (min) & (\%) \\
\midrule
AA & Gen & 12.89 & -34.0 & 3298.0 & 3.73 \\
AA & Feb & 13.28 & -27.0 & 3403.0 & 1.32 \\
AA & Mar & 15.85 & -27.0 & 3288.0 & 1.52 \\
AA & Apr & 15.96 & -28.0 & 2595.0 & 1.43 \\
\midrule
DL & Gen & 12.83 & -30.0 & 1209.0 & 2.80 \\
DL & Feb & 11.35 & -30.0 & 1222.0 & 0.23 \\
DL & Mar & 10.32 & -32.0 & 1176.0 & 0.49 \\
DL & Apr &  8.71 & -30.0 & 1145.0 & 0.51 \\
\bottomrule
\end{tabular}
\end{table}

\textbf{Analisi}: entrambe le compagnie mostrano ritardi medi più alti in
gennaio, periodo caratterizzato da condizioni meteo avverse. AA presenta
ritardi medi costantemente superiori a DL (fino a +7 min ad aprile) e un tasso
di cancellazione significativamente più alto a gennaio (3.73\% vs 2.80\%).
Delta mostra un miglioramento progressivo del ritardo medio nel corso dei
quattro mesi, mentre AA rimane sostanzialmente stabile tra febbraio e aprile.
I valori massimi anomali di AA (fino a 3403 min) suggeriscono la presenza di
outlier estremi.

\subsection{Q2 — Top-10 compagnie per ritardo medio in arrivo}

La query Q2 considera i voli non cancellati e non deviati
(\texttt{CANCELLED = 0} e \texttt{DIVERTED = 0}) con almeno 500 voli nel
periodo, e calcola il ranking delle 10 compagnie con il maggiore ritardo medio
in arrivo (\texttt{ARR\_DELAY}). I valori \texttt{NULL} nelle componenti di
ritardo sono trattati come zero, in accordo con la convenzione BTS (i campi
sono nulli quando il ritardo totale è inferiore a 15 minuti).

I risultati sono in Tabella~\ref{tab:q2}.

\begin{table}[h]
\caption{Q2 — Top-10 compagnie per ritardo medio in arrivo.}
\label{tab:q2}
\small
\begin{tabular}{lrrrrrrr}
\toprule
Carrier & Voli & Arr & Carr & Wthr & NAS & Sec & Late \\
\midrule
OH &  78.476 & 17.16 & 7.24 & 1.66 & 3.31 & 0.03 & 12.13 \\
F9 &  65.965 & 11.32 & 5.92 & 0.49 & 3.95 & 0.00 & 10.42 \\
G4 &  42.596 &  9.54 & 5.68 & 2.04 & 4.27 & 0.16 &  7.54 \\
AA & 302.790 &  8.98 & 5.95 & 1.01 & 2.97 & 0.02 &  7.41 \\
OO & 259.692 &  7.59 & 7.93 & 2.47 & 2.68 & 0.02 &  3.38 \\
B6 &  75.680 &  6.15 & 4.49 & 0.43 & 4.31 & 0.03 &  7.08 \\
HA &  25.707 &  5.81 & 4.46 & 0.54 & 0.10 & 0.02 &  3.29 \\
MQ &  87.416 &  5.32 & 3.09 & 1.55 & 3.06 & 0.02 &  5.15 \\
DL & 312.503 &  4.19 & 5.87 & 0.89 & 2.79 & 0.01 &  4.06 \\
UA & 248.433 &  2.46 & 3.39 & 0.61 & 3.83 & 0.00 &  4.48 \\
\bottomrule
\end{tabular}
\end{table}

\textbf{Analisi}: Endeavor Air (OH) risulta la compagnia con il ritardo medio
più alto (17.16 min), con la componente ``Late Aircraft'' come principale
contributore (12.13 min), indicando una forte propagazione dei ritardi tra
voli consecutivi. Frontier (F9) mostra un profilo simile. OO (SkyWest) ha
invece la componente ``Carrier'' come dominante (7.93 min), suggerendo
problemi interni operativi. United (UA) è la compagnia meglio performante
nella top-10, con solo 2.46 min di ritardo medio in arrivo.

\subsection{Q3 — Percentili di ritardo per fascia oraria}

La query Q3 calcola, per AA, DL, UA e WN, i percentili p25, p50, p75 e p90 di
\texttt{DEP\_DELAY} per ciascuna delle 24 fasce orarie, considerando solo voli
non cancellati. I percentili sono calcolati con la funzione
\texttt{percentile\_approx()} di Spark, che implementa un algoritmo efficiente
per grandi volumi di dati. I valori globali di min/max per carrier sono
riportati in Tabella~\ref{tab:q3range}.

\begin{table}[h]
\caption{Q3 — Min e Max di DEP\_DELAY per carrier.}
\label{tab:q3range}
\begin{tabular}{lrr}
\toprule
Carrier & Min (min) & Max (min) \\
\midrule
AA & -34.0 & 3403.0 \\
DL & -32.0 & 1222.0 \\
UA & -50.0 & 1349.0 \\
WN & -24.0 &  721.0 \\
\bottomrule
\end{tabular}
\end{table}

\textbf{Analisi}: tutte le compagnie mostrano un pattern comune: i voli
mattutini (ore 5--8) hanno ritardi mediani negativi o prossimi a zero (anticipo
medio), mentre i voli serali (ore 17--21) presentano ritardi sistematicamente
più alti, con p90 che supera i 50--70 minuti. Questo effetto si spiega con
l'accumulo di ritardi nel corso della giornata (chain delay). Southwest (WN)
mostra ritardi assoluti più contenuti rispetto alle altre compagnie, con un
massimo di 721 minuti contro i 3403 di AA. United presenta il valore minimo
più basso (-50 min), indicando sistematicamente partenze in anticipo.

% ─────────────────────────────────────────────────────────────────────────────
\section{Analisi di Clustering}

\subsection{Feature Engineering}

Per ciascuna delle 15 compagnie con il maggior numero di voli nel periodo, è
stato calcolato un vettore di 12 feature aggregate:

\begin{enumerate}
  \item \texttt{avg\_dep\_delay}: ritardo medio in partenza (voli non cancellati)
  \item \texttt{avg\_arr\_delay}: ritardo medio in arrivo (voli completati)
  \item \texttt{cancellation\_rate}: percentuale di voli cancellati
  \item \texttt{avg\_carrier\_delay}, \texttt{avg\_weather\_delay},
        \texttt{avg\_nas\_delay}, \texttt{avg\_security\_delay},
        \texttt{avg\_late\_aircraft}: medie delle 5 componenti di ritardo
  \item \texttt{std\_dep\_delay}: deviazione standard del ritardo in partenza
  \item \texttt{on\_time\_rate}: percentuale di voli con \texttt{DEP\_DELAY} $\leq$ 0
  \item \texttt{diverted\_rate}: percentuale di voli deviati
  \item \texttt{avg\_delay\_recovery}: media di (\texttt{DEP\_DELAY} $-$ \texttt{ARR\_DELAY})
\end{enumerate}

Le feature sono state normalizzate tramite \texttt{StandardScaler} (media zero,
deviazione standard unitaria) prima dell'applicazione del clustering.

\subsection{Selezione del Numero di Cluster}

Il valore ottimale di $k$ è stato determinato combinando due metriche:
l'Elbow Method (basato su WSSSE) e l'indice di Silhouette, valutati per
$k \in \{2, 3, 4, 5, 6\}$. Il valore massimo del Silhouette score è stato
ottenuto per $k = 2$, confermando la presenza di due gruppi nettamente
distinti nel dataset.

\subsection{Risultati e Interpretazione}

Il clustering con $k = 2$ ha prodotto i seguenti gruppi:

\textbf{Cluster 0 — ``Carrier a prestazioni medie''} (11 carrier: AS, B6, DL,
G4, HA, MQ, NK, OO, UA, WN, YX): compagnie con ritardi contenuti e tassi di
cancellazione moderati. Delta (DL) e Alaska (AS) rappresentano le eccellenze
operative del gruppo.

\textbf{Cluster 1 — ``Carrier ad alto ritardo''} (3 carrier: AA, F9, OH):
compagnie caratterizzate da ritardi medi in partenza superiori a 14 minuti,
forte componente \texttt{Late Aircraft} e valori elevati di
\texttt{std\_dep\_delay} (alta imprevedibilità). Endeavor Air (OH) è il caso
estremo con 19.79 min di avg dep delay.

La proiezione PCA a 2 componenti mostra una separazione netta tra i due
cluster, con il Cluster 1 isolato nella zona di alto ritardo. Le prime due
componenti principali spiegano la maggior parte della varianza del dataset,
confermando la qualità della separazione.

% ─────────────────────────────────────────────────────────────────────────────
\section{Valutazione delle Prestazioni}

\subsection{Piattaforma di Riferimento}

I job sono stati eseguiti su un nodo standalone in modalità Docker Compose
(Windows 11, WSL2, modalità \texttt{local[2]}). La configurazione Spark
prevede 1 worker con 1536 MB di memoria e 2 core.

\subsection{Confronto DataFrame API vs Spark SQL}

Le tre query sono state implementate con entrambe le API (DataFrame e Spark
SQL) e confrontate in termini di tempo di esecuzione wall-clock. I tempi sono
stati misurati dallo script \texttt{compare\_queries.py}, che esegue le due
versioni in sequenza sullo stesso dataset cachato e verifica la correttezza
dei risultati tramite confronto bidirezionale (\texttt{subtract}).

\begin{table}[h]
\caption{Tempi di esecuzione (secondi) su nodo standalone.}
\label{tab:timing}
\begin{tabular}{lrrr}
\toprule
Query & DataFrame API (s) & Spark SQL (s) & $\Delta$ (s) \\
\midrule
Q1 & TODO & TODO & TODO \\
Q2 & TODO & TODO & TODO \\
Q3 & TODO & TODO & TODO \\
\bottomrule
\end{tabular}
\end{table}

% ── Istruzioni per compilare la tabella ────────────────────────────────────
% Eseguire: docker exec spark-master python3 /opt/spark/jobs/compare_queries.py
% I tempi vengono stampati nel riepilogo finale. Sostituire i TODO sopra.
% ───────────────────────────────────────────────────────────────────────────

\textbf{Analisi}: % [da completare dopo aver eseguito compare_queries.py]
In generale, Spark SQL e DataFrame API producono risultati identici (differenze
pari a zero righe in entrambe le direzioni). Le differenze di tempo sono
attribuibili principalmente all'overhead del parser SQL e all'ottimizzazione
del Catalyst optimizer, che in entrambi i casi genera lo stesso piano fisico
di esecuzione.

% ─────────────────────────────────────────────────────────────────────────────
\section{Conclusioni}

Il sistema implementato dimostra l'efficacia dell'integrazione di tecnologie
open-source per l'analisi Big Data: NiFi garantisce un'ingestion robusta e
configurabile, Spark permette un processing efficiente sia tramite DataFrame
API che Spark SQL, e la combinazione Redis+Grafana offre una visualizzazione
interattiva dei risultati.

I risultati analitici evidenziano differenze operative significative tra le
compagnie aeree: AA mostra ritardi sistematicamente più alti di DL in tutti i
mesi; OH (Endeavor) è la compagnia con il peggior ritardo medio in arrivo,
trainato dalla componente Late Aircraft; il pattern di accumulo dei ritardi nel
corso della giornata è trasversale a tutte le compagnie. Il clustering ha
identificato due cluster netti: un gruppo di carrier con prestazioni nella
media e un gruppo ristretto (AA, F9, OH) caratterizzato da ritardi
strutturalmente elevati.

\bibliographystyle{ACM-Reference-Format}
\begin{thebibliography}{9}

\bibitem{dunning2021}
T. Dunning.
\newblock The t-digest: Efficient estimates of distributions.
\newblock \emph{Software Impacts}, 7:100049, 2021.

\bibitem{grafana}
Grafana Labs.
\newblock Grafana.
\newblock \url{https://grafana.com/grafana/}.

\bibitem{karnin2016}
Z. Karnin, K. Lang, and E. Liberty.
\newblock Optimal quantile approximation in streams.
\newblock In \emph{Proc. of IEEE 57th FOCS}, pages 71--78, 2016.

\end{thebibliography}

\end{document}
